\centerline{\bf \Huge Окружность Аполлония}

\begin{flushright}
        \scriptsize{Эпизод 15. Ложь и молчание.}
\end{flushright}

\problem{1}
\textbf{(Окружность Аполлония).} Даны точки $A$ и $B$ и положительное число $k$. Докажите, что при $k \neq 1$ множество точек $X$, удовлетворяющих условию $A X: B X=k$ это окружность, центр которой лежит на прямой $A B$.

\problem{2}
Обозначим за $\Omega_a$ окружность, построенную как на диаметре на отрезке, соединяющем основания внешней и внутренней биссектрис треугольника $A B C$, проведенных из вершины $A$. Окружности $\Omega_b$ и $\Omega_c$ определяются аналогично.

(a) Докажите, что три построенные окружности имеют две общие точки, которые называются двумя \textit{точками Аполлония.}

(b) Точка $A_1$ на прямой $B C$ такова, что прямая $A_1 A$ касается окружности $(A B C)$. Точки $B_1$ и $C_1$ определяются аналогично. Выведите из предыдущего пункта, что точки $A_1, B_1$ и $C_1$ лежат на одной прямой.

\problem{3}
Докажите, что на прямой, соединяющей две точки Аполлония треугольника лежит:

(a) Центр описанной около этого треугольника окружности.

(b) Точка Лемуана этого треугольника.

(c) Ортоцентр треугольника, вершины которого являются основаниями внутренних биссектрис треугольника.


\problem{4}
(a) Докажите, что точки Аполлония изогонально сопряжены двум точкам Торричелли треугольника.
(b) Докажите, что проекции точки Аполлония на стороны треугольника являются вершинами равностороннего треугольника.

\problem{5}
Точка $P$ внутри неравнобедренного треугольника $A B C$ такова, что существует окружность, проходящая через точки $A$ и $P$, касающаяся окружностей ( $A B C$ ) и ( $B P C$ ). Докажите, что биссектрисы углов $B P C$ и $B A C$ пересекаются на стороне $B C$.

\newpage

\hspace{3mm}

\problem{6}
\textbf{Теорема.} В треугольнике $A B C$ проведена тройка чевиан $A A_1, B B_1$ и $C C_1$, пересекающиеся в одной точке. Прямая $B_1 C_1$ пересекает $B C$ в точке $A_2$, точки $B_2$ и $C_2$ определяются аналогично. Обозначим за $\Omega_a$ окружность, построенную на отрезке $A_1 A_2$ как на диаметре, аналогично определим окружности $\Omega_b$ и $\Omega_c$.

(a) Докажите, что степень центра описанной окружности треугольника $(A B C)$ одна и та же относительно всех трех построенных окружностей.

(b) Докажите, что степень ортоцентра треугольника $A_1 B_1 C_1$ одинакова относительно всех трех построенных окружностей.

(c) Докажите, что середины отрезков $A_1 A_2, B_1 B_2$ и $C_1 C_2$ лежат на одной прямой $\ell$.

(d) Докажите, что прямая, соединяющая ортоцентр треугольника $A_1 B_1 C_1$ с центром описанной окружности треугольника $A B C$ перпендикулярна прямой $\ell$.

\problem{7}
Дан треугольник $A B C$ и такая точка $P$, что $A P: B C=B P: C A=C P: A B$. Докажите, что точка $P$ лежит на прямой Эйлера треугольника $A B C$.

\problem{8}
Отметим середины трех биссектрис треугольника. Докажите, что ортоцентр треугольника из этих середин лежит на прямой Эйлера исходного треугольника.